% !PRINTED 33
\weberpage{33}

For divisibility of functions the same laws hold as for divisibility of numbers, in particular the following:

\gesperrt{1. If the function $F(x)$ is divisible by the function $f(x)$, and $f(x)$ is divisible by a third function $\varphi(x)$, then $F(x)$ is also divisible by $\varphi(x)$.}

For if
\[
 F=Qf,
 \qquad
 f=q\varphi,
\]
where $Q$ and $q$ are whole rational functions, then
\[
 F=Qq\varphi,
\]
and since $Qq$ is a whole rational function, $F$ is divisible by $\varphi$.

\gesperrt{2. If $F(x)$ is divisible by $f(x)$ and $Q$ is any whole rational function, then $QF(x)$ is also divisible by $f(x)$.}

\gesperrt{3. If $F(x)$ and $f(x)$ are divisible by $\varphi(x)$, then $F(x)\pm f(x)$ is also divisible by $\varphi(x)$, or more generally:}

\gesperrt{4. If $F_1,F_2\ldots$ are divisible by $f(x)$ and $Q_1,Q_2\ldots$ are arbitrary whole rational functions, then}
\[
 Q_1F_1+Q_2F_2+\cdots
\]
\gesperrt{is also divisible by $f(x)$.}

The last theorem includes the two preceding ones and is proved simply as follows.

If $F_1,F_2\ldots$ are divisible by $f$, then the whole rational functions $\Phi_1,\Phi_2\ldots$ can be determined so that
\[
 F_1=\Phi_1f,
 \qquad
 F_2=\Phi_2f,\ldots
\]
and consequently
\[
 Q_1F_1+Q_2F_2+\cdots=(Q_1\Phi_1+Q_2\Phi_2+\cdots)f.
\]
Since $Q_1\Phi_1+Q_2\Phi_2+\cdots$ is a whole rational function, the theorem is proved.

\gesperrt{5. Every function is divisible by itself.}

With respect to divisibility or indivisibility of functions nothing is changed if the functions are multiplied by arbitrary factors independent of $x$.

A non-zero quantity independent of $x$ may be regarded as a function of degree zero. If we call such a quantity a constant, we may say:

\gesperrt{6. Every function is divisible by every constant.}
